\chapter{Evaluación de Plataformas}

La elección del stack tecnológico se realizó comparando las alternativas
disponibles en tres dimensiones: framework de desarrollo móvil, backend como
servicio y estrategia de almacenamiento y notificaciones.

\section{Framework de desarrollo móvil}

Se evaluaron Flutter, Kotlin nativo (Android) y React Native. Kotlin nativo
garantiza el máximo desempeño y acceso directo a las APIs de Android, pero
restringe la solución a una única plataforma y aumenta el esfuerzo de desarrollo
para una segunda plataforma. React Native permite compartir código entre
plataformas, pero introduce dependencias de puentes nativos y mayor variabilidad
en el desempeño. Flutter se selecciona porque permite construir interfaces
consistentes en Android e iOS desde una única base de código en Dart, ofrece un
sistema de widgets con Material 3 que acelera el desarrollo de interfaces
accesibles y cuenta con un ecosistema maduro de paquetes (Riverpod, go\_router,
cryptography). Para un proyecto académico que debe demostrar resultados en
plazos acotados, Flutter presenta el mejor equilibrio entre velocidad de
desarrollo, calidad visual y soporte.

\begin{table}[H]
\centering
\small
\tablazebra
\caption{Comparación de frameworks de desarrollo móvil.}
\label{tab:comparacion-frameworks}
\begin{tabularx}{\textwidth}{P{0.78}P{1.07}P{1.07}P{1.08}}
\toprule
\filaencabezado \encab{Criterio} & \encab{Flutter} & \encab{Kotlin nativo} & \encab{React Native} \\
\midrule
Plataformas &
Android e iOS con una base de código &
Solo Android &
Android e iOS con puentes nativos \\
Desempeño &
Alto (compila a nativo) &
Máximo &
Variable (depende del puente) \\
Ecosistema de paquetes &
Maduro (Riverpod, go\_router, cryptography) &
Maduro para Android &
Amplio, con mayor variabilidad \\
Velocidad de desarrollo &
Alta (Material 3, widgets) &
Media &
Media \\
\bottomrule
\end{tabularx}
\end{table}

\section{Backend como servicio}

Se compararon Firebase, Supabase y un backend propio. Un backend propio ofrece
control total, pero exige mantener servidores, autenticación, seguridad y
despliegue, lo que excede el alcance y las horas disponibles. Supabase, basado
en PostgreSQL, es una alternativa válida de open source, pero requiere
configuración adicional de políticas de seguridad y su ecosistema de integración
móvil es menor. Firebase se selecciona por su integración nativa con Flutter, la
sincronización en tiempo real de Cloud Firestore, la persistencia local
integrada que habilita el modo offline-first y la disponibilidad del plan
gratuito (Spark) suficiente para un piloto académico.

\section{Almacenamiento local}

Se evaluó el uso exclusivo de Room sobre SQLite (opción considerada en el
anteproyecto) frente a Firestore. Room ofrece almacenamiento 100~\% local y
máximo control, pero dificulta la sincronización y la persistencia en múltiples
dispositivos. Firestore, si bien es un servicio en la nube, compensa la
dependencia de red con tres mecanismos: persistencia local habilitada, caché de
metadatos en SharedPreferences y caché de archivos con auto-descarga. Esta
combinación entrega un comportamiento equivalente al offline esperado por el
anteproyecto, sumando el beneficio de la sincronización entre el dispositivo del
cuidador y los servicios de respaldo.

\section{Notificaciones}

Se evaluaron notificaciones push con FCM frente a notificaciones locales. FCM
permite alertar al usuario incluso con la aplicación cerrada y desde el
servidor, pero exige infraestructura adicional, manejo de permisos de plataforma
y puede fallar sin conexión. Para los recordatorios de valoración diaria del
paciente se opta por notificaciones locales (flutter\_local\_notifications),
re-planificadas al iniciar sesión: no dependen de red, respetan horas exactas
mediante SCHEDULE\_EXACT\_ALARM y cubren completamente el requerimiento
de recordar al cuidador la valoración diaria.

\section{Conclusión de la evaluación}

Como conclusión, el stack elegido —Flutter/Dart con Riverpod y go\_router,
Firebase (Auth, Firestore, Functions, Storage) y notificaciones locales—
satisface los requerimientos funcionales y no funcionales del proyecto, es de
costo cero, está ampliamente documentado y permite un funcionamiento offline
-first con seguridad por diseño.